% Part: first-order-logic
% Chapter: completeness
% Section: outline

\documentclass[../../../include/open-logic-section]{subfiles}

\begin{document}

\iftag{FOL}
      {\olfileid{fol}{com}{out}}
      {\olfileid{pl}{com}{out}}

\olsection{Overzicht van het bewijs}

Het bewijs van de volledigheidsstelling is enigszins complex en bij
een eerste lezing raakt men gemakkelijk de draad kwijt. Daarom geven
we eerst een overzicht. De eerste stap is een verandering van
gezichtspunt waardoor een route naar het bewijs zichtbaar wordt. Vat
men volledigheid op als ``als $\Gamma \Entails !A$, dan $\Gamma
\Proves !A$'', dan is het moeilijk zelfs maar een idee voor het bewijs
te vinden. Om $\Gamma \Proves !A$ aan te tonen, moeten we immers
!!a{derivation} vinden, terwijl de aanname $\Gamma \Entails !A$ ons
daarbij op het eerste gezicht niet helpt. Bij sommige bewijssystemen
kan men rechtstreeks !!a{derivation} construeren, maar wij kiezen een
iets andere benadering. De vereiste verandering van gezichtspunt is
deze: volledigheid kan ook worden geformuleerd als ``als $\Gamma$
consistent is, dan is zij vervulbaar.'' Misschien kunnen we de
informatie in~$\Gamma$ en de aanname dat zij consistent is gebruiken
om \iftag{FOL}{!!a{structure}}{!!a{valuation}} te construeren die elke
\iftag{FOL}{!!{sentence}}{!!{formula}} in~$\Gamma$ vervult. We weten
tenslotte wat voor \iftag{FOL}{!!{structure}}{!!{valuation}} we zoeken:
een die precies is zoals $\Gamma$ haar beschrijft!{}

Als $\Gamma$ alleen \iftag{FOL}{atomaire
  !!{sentence}s}{!!{propositional variable}s} bevat, is het eenvoudig
er een model voor te construeren.\iftag{FOL}{ Veronderstel dat alle
  atomaire !!{sentence}s de vorm $\Atom{P}{a_1,\dots,a_n}$ hebben,
  waarbij de $a_i$ !!{constant}s zijn.}{} We hoeven slechts
\iftag{FOL}{%
  !!a{domain}~$\Domain{M}$ en een interpretatie voor~$P$ te kiezen
  zodat $\Sat{M}{\Atom{P}{a_1,\ldots,a_n}}$. Dat is niet moeilijk: stel
  $\Domain{M} = \Nat$ en $\Assign{\Obj c_i}{M} = i$, en neem voor elke
  $\Atom{P}{a_1,\ldots,a_n} \in \Gamma$ het tupel $\tuple{k_1,
    \dots, k_n}$ op in $\Assign{P}{M}$, waarbij $k_i$ de index is van
  het constante symbool~$a_i$ (dus $a_i \ident \Obj c_{k_i}$).}{%
  !!a{valuation}~$\pAssign{v}$ te kiezen waarvoor $\pSat{v}{p}$ voor
  alle $p \in \Gamma$. Stel daartoe $\pAssign{v}(p) = \True$ dan en
  slechts dan als $p \in \Gamma$.}

Veronderstel nu dat $\Gamma$ een !!{formula}~$\lnot !B$ bevat, waarbij
$!B$ atomair is. Men zou kunnen vrezen dat de constructie van
$\iftag{FOL}{\Struct{M}}{\pAssign{v}}$ verhindert dat $\lnot !B$ waar
wordt. Hier komt de consistentie van~$\Gamma$ echter van pas: als
$\lnot !B \in \Gamma$, dan $!B \notin \Gamma$, want anders zou
$\Gamma$ inconsistent zijn. Als $!B \notin \Gamma$, dan geldt volgens
onze constructie van~$\iftag{FOL}{\Struct{M}}{\pAssign{v}}$ dat
$\iftag{FOL}{\Sat/{M}}{\pSat/{v}}{!B}$, en dus
$\iftag{FOL}{\Sat{M}}{\pSat{v}}{\lnot !B}$. Tot zover gaat alles goed.

Wat gebeurt er als $\Gamma$ samengestelde, niet-atomaire formules
bevat? Stel dat zij $!A \land !B$ bevat. Om die formule waar te maken,
moeten we te werk gaan alsof zowel $!A$ als $!B$ tot~$\Gamma$ behoort.
En als $!A \lor !B \in \Gamma$, moeten we ten minste een van beide
waar maken, dus te werk gaan alsof een van beide tot~$\Gamma$ behoort.

Dit leidt tot het volgende idee. We voegen !!{formula}s aan~$\Gamma$
toe op zo'n manier dat (a)~de resulterende verzameling consistent
blijft; (b)~voor elke mogelijke atomaire !!{sentence}~$!A$ ofwel $!A$
ofwel $\lnot !A$ in de resulterende verzameling staat; en (c)~als
$!A \land !B$ in de verzameling staat, zowel $!A$ als $!B$ erin staat,
als $!A \lor !B$ erin staat, ten minste een van $!A$ en $!B$ erin staat,
enzovoort. We zetten dit proces voort, mogelijk onbeperkt. Noem de
verzameling van alle zo toegevoegde !!{formula}s~$\Gamma^*$. Onze
bovenstaande constructie zou dan
\iftag{FOL}{!!a{structure}~$\Struct{M}$}{!!a{valuation}~$\pAssign{v}$}
opleveren waarvan we door inductie kunnen bewijzen dat zij alle zinnen
in~$\Gamma^*$ vervult, en dus ook alle zinnen in~$\Gamma$, omdat
$\Gamma \subseteq \Gamma^*$. Het blijkt voldoende om (a) en~(b) te
garanderen. Een verzameling zinnen die aan (b) voldoet, heet
\emph{volledig}. Onze taak is dus de consistente verzameling~$\Gamma$
uit te breiden tot een consistente en volledige verzameling~$\Gamma^*$.

\iftag{FOL}{%
Dit plan heeft één moeilijkheid. Als $\lexists[x][!A(x)] \in \Gamma$,
zouden we graag een !!{constant}~$c$ kiezen en tijdens het proces
$!A(c)$ toevoegen. Maar hoe weten we dat dit altijd mogelijk is?
Misschien bevat onze taal slechts enkele !!{constant}s en geldt voor
elk daarvan $\lnot !A(c) \in \Gamma$. We kunnen dan niet ook $!A(c)$
toevoegen, want daardoor zou de verzameling inconsistent worden en
zouden we niet weten of $\Struct{M}$ de formule $!A(c)$ dan wel
$\lnot !A(c)$ waar moet maken. Het kan bovendien voorkomen dat
$\Gamma$ alleen zinnen bevat in een taal zonder enig constant symbool,
zoals de taal van de verzamelingenleer.

De oplossing is aan het begin eenvoudigweg oneindig veel constanten
toe te voegen, samen met zinnen die deze op de juiste manier met de
kwantoren verbinden. Uiteraard moeten we nagaan dat hierdoor geen
inconsistentie kan ontstaan.

Onze oorspronkelijke constructie werkt goed als de atomaire zinnen
alleen !!{constant}s bevatten. De taal kan echter ook !!{function}s
bevatten. Dan kan het lastig zijn om aan deze !!{function}s geschikte
functies op~$\Nat$ toe te kennen waarmee de constructie werkt. We
gebruiken daarom een andere kunstgreep. In plaats van $i$ te gebruiken
als interpretatie van $\Obj c_i$, nemen we de verzameling
!!{constant}s zelf als domein. Dan kan $\Struct M$ elke !!{constant}
aan zichzelf toekennen: $\Assign{\Obj c_i}{M} = \Obj c_i$. We kunnen
nog verder gaan en $\Domain{M}$ gelijkstellen aan de verzameling van
alle \emph{termen} van de taal!{} Er is dan een voor de hand liggende
toekenning aan de !!{function}s van functies die termen als argumenten
en als waarden hebben: aan de !!{function}~$\Obj f^n_i$ kennen we de
functie toe die bij $n$ invoertermen $t_1$, \dots,~$t_n$ de term
$\Obj f^n_i(t_1, \dots, t_n)$ als waarde oplevert.

Het laatste stuk van de puzzel is de behandeling van~$\eq$. Het
!!{predicate}~$\eq$ heeft een vaste interpretatie:
$\Sat{M}{\eq[t][t']}$ dan en slechts dan als
$\Value{t}{M} = \Value{t'}{M}$. Als we de constructie zo inrichten dat
de !!{value} van een term~$t$ de term $t$ zelf is, maakt deze
!!{structure} dus \emph{geen enkele} zin van de vorm $\eq[t][t']$ waar,
tenzij $t$ en $t'$ precies dezelfde term zijn. Dat vormt een probleem,
want vrijwel elke belangwekkende theorie in een taal met !!{function}s
heeft zinnen $\eq[t][t']$ als stelling waarbij $t$ en $t'$ niet dezelfde
term zijn, bijvoorbeeld in rekenkundige theorieën:
$\eq[(\Obj 0+ \Obj 0)][\Obj 0]$. Om dit probleem op te lossen,
veranderen we het domein van~$\Struct M$. In plaats van termen als de
objecten van~$\Domain{M}$ te gebruiken, gebruiken we verzamelingen van
termen. Elke zo'n verzameling bevat alle termen waarvan de zinnen
in~$\Gamma$ vereisen dat ze gelijk zijn. Als $\Gamma$ bijvoorbeeld een
rekenkundige theorie is, zal een van deze verzamelingen $\Obj 0$,
$(\Obj 0 + \Obj 0)$, $(\Obj 0 \times \Obj 0)$ enzovoort bevatten. Dit
is de verzameling die we aan $\Obj 0$ toekennen, en zij blijkt ook de
waarde te zijn van alle termen die zij bevat, dus bijvoorbeeld ook van
$(\Obj 0 + \Obj 0)$. De zin $\eq[(\Obj 0+ \Obj 0)][\Obj 0]$ is daarom
waar in deze aangepaste !!{structure}.}{}

We gaan dus als volgt te werk. Eerst onderzoeken we de eigenschappen
van consistente, !!{complete} verzamelingen. In het bijzonder bewijzen
we dat !!a{complete} consistente verzameling $!A \land !B$ bevat dan en
slechts dan als zij zowel $!A$ als~$!B$ bevat, en $!A \lor !B$ bevat
dan en slechts dan als zij ten minste een van beide bevat, enzovoort
(\olref[ccs]{prop:ccs}).\iftag{FOL}{ Vervolgens definiëren en
  onderzoeken we ``verzadigde'' verzamelingen zinnen. Een verzadigde
  verzameling bevat implicaties die elke gekwantificeerde !!{sentence}
  met haar instanties verbinden (\olref[hen]{defn:henkin-exp}). We
  bewijzen dat elke consistente verzameling~$\Gamma$ kan worden
  uitgebreid tot een verzadigde verzameling~$\Gamma'$
  (\olref[hen]{lem:henkin}). Als een verzameling consistent, verzadigd
  en !!{complete} is, heeft zij bovendien de eigenschap dat zij
  \iftag{prvEx}{$\lexists[x][!A(x)]$ bevat dan en slechts dan als zij
    $!A(t)$ bevat voor een gesloten term~$t$}{}\iftag{defEx,defAll}{}{ en
  }\iftag{prvAll}{$\lforall[x][!A(x)]$ bevat dan en slechts dan als zij
    $!A(t)$ bevat voor alle gesloten termen~$t$}{}
  (\olref[hen]{prop:saturated-instances}).}{}
Daarna nemen we de\iftag{FOL}{ verzadigde}{} consistente
verzameling~$\iftag{FOL}{\Gamma'}{\Gamma}$ en tonen we aan dat zij kan
worden uitgebreid tot een \iftag{FOL}{verzadigde, consistente en}{consistente en}
!!{complete} verzameling~$\Gamma^*$ (\olref[lin]{lem:lindenbaum}). Met
deze verzameling $\Gamma^*$ definiëren we ons \iftag{FOL}{termmodel
  $\Struct M(\Gamma^*)$}{!!{valuation}~$\pAssign v(\Gamma^*)$}.
\iftag{FOL}{Het termmodel heeft de verzameling gesloten termen als
  domein, en de interpretatie van zijn !!{predicate}s wordt gegeven
  door de atomaire !!{sentence}s}{De valuatie wordt bepaald door de
  !!{propositional variable}s} in~$\Gamma^*$
(\olref[mod]{defn:termmodel}). We
gebruiken de eigenschappen van\iftag{FOL}{ verzadigde,}{} volledige
consistente verzamelingen om aan te tonen dat inderdaad
$\iftag{FOL}{\Sat{M(\Gamma^*)}}{\pSat{v(\Gamma^*)}}{!A}$ dan en slechts
dan als $!A \in \Gamma^*$ (\olref[mod]{lem:truth}), en dus in het
bijzonder $\iftag{FOL}{\Sat{M(\Gamma^*)}}{\pSat{v(\Gamma^*)}}{\Gamma}$.
\iftag{FOL}{Ten slotte onderzoeken we hoe een termmodel moet worden
  gedefinieerd als $\Gamma$ ook~$\eq$ bevat
  (\olref[ide]{defn:term-model-factor}) en bewijzen we dat dit
  model~$\Gamma^*$ vervult (\olref[ide]{lem:truth}).}{}

\end{document}
